\section{Performance Investigation}
\label{pi}

\begin{table}
%\small
  \centering
  \setlength{\tabcolsep}{5mm}
%  \renewcommand{\arraystretch}{0.8}
\begin{tabular}{cc}
\toprule
\multicolumn{1}{c}{Edge Pruning Strategy}     & RA $\uparrow$ \\ \midrule
Top-3     & 54.00         \\
Bot-3                & 56.00         \\
AMT       & 52.00         \\
AAT (Ours) & \textbf{57.00}     \\  \bottomrule
\end{tabular}
  \caption{\label{average_based}
   Task performance of CortexDebate with different edge pruning strategies. ``$\uparrow$'' means that higher values are better. The best record is highlighted in bold.
  }
\end{table}
\begin{table}
%\small
  \centering
  \setlength{\tabcolsep}{3mm}
%  \renewcommand{\arraystretch}{0.8}
\begin{tabular}{cc}
\toprule
\multicolumn{1}{c}{Text Similarity Calculation}     & M-Avg $\uparrow$ \\ \midrule
DeTS     & 57.57{\tiny $\pm$0.69}         \\
R1-P                & 56.91{\tiny $\pm$3.10}         \\
Ours & \textbf{60.31{\tiny $\pm$0.32}}     \\  \bottomrule
\end{tabular}
  \caption{\label{text_similarity}
   Task performance of CortexDebate with different calculation strategies of text similarity. The format of the results is ``(average result)$\pm$(variance)''. ``$\uparrow$'' means that higher values are better. The best record is highlighted in bold.
  }
\end{table}

In this section, we present additional in-depth investigations on our CortexDebate to further analyze its effectiveness.


\subsection{Average-based Edge Pruning}

As present in Equation~\eqref{11}, our CortexDebate uses the average value as the threshold to prune the edges. To evaluate the effectiveness of adopting the average-based threshold, we conduct experiments on MATH dataset. Specifically, we compare the strategy of debating with the agents above the average threshold (AAT) with three alternatives, namely Top-3 (debating with the top three agents), Bot-3 (excluding the bottom three agents), and AMT (debating with agents above the median threshold). As shown in Table~\ref{average_based}, AAT achieves superior performance than the others, which validates the effectiveness of the average-based edge pruning.


\subsection{Text Similarity Calculation}

In the experiments, we use the text-embedding-3-large model released by OpenAI as the embedding model. For calculation of text similarity, we use cosine similarity which is widely used due to its efficiency and robustness to varying input length, and strong empirical performance. To confirm this, on LongBench dataset, we compare our strategy with ag-nli-DeTS-sentence-similarity-v4 model (DeTS) that is trained on six natural language inference datasets and “DeepSeek R1 + prompt” (R1-P). The average task scores and variances over three runs are reported in Table~\ref{text_similarity}. We can find that our method outperforms the two baseline methods. Moreover, we observe that large language models tend to perform unstably when used for text similarity estimation, which is reflected by the higher variance across different runs.


\subsection{Confidence Recalibration}

As present in Equation~\eqref{recalibrate}, our CortexDebate adopts a recalibration strategy to mitigate overconfidence issue of LLM agents. Following the practice of~\cite{chen2023reconcile}, we conduct experiments on different combinations of thresholds. We present the average scores on eight adopted datasets in Table~\ref{confidence}. According to the results, this recalibration strategy leads to the improved performance, and the parameter setting in our paper leads to the best performance. Moreover, our method is generally insensitive to different threshold configurations. This demonstrates the effectiveness and robustness of the proposed confidence recalibration strategy.

\begin{table}
%\small
  \centering
  \setlength{\tabcolsep}{2mm}
%  \renewcommand{\arraystretch}{0.8}
\begin{tabular}{cc}
\toprule
Threshold Configurations     & Average Score $\uparrow$ \\ \midrule
w/o recalibration  & 68.62    \\
$[$0.9, 0.7, 0.2$]$    & 68.96    \\
$[$0.9, 0.6, 0.3$]$    & 69.03    \\
$[$0.8, 0.6, 0.2$]$    & 69.15    \\
$[$0.8, 0.5, 0.1$]$    & 68.84    \\
$[$0.7, 0.5, 0.2$]$    & 68.72    \\
$[$0.8, 0.6, 0.3$]$ (Ours)    & \textbf{69.41}    \\
  \bottomrule
\end{tabular}
  \caption{\label{confidence}
   Task performance of CortexDebate with different threshold configurations of the recalibration strategy. ``$\uparrow$'' means that higher values are better. The best record is highlighted in bold.
  }
\end{table}